\section*{Chapter \ref{ch:b2:endogenous-growth} --- Ideas and Endogenous Growth}

\begin{solution}{exo:b2:endogenous-growth:1}
% figdata: tfp, res_summary
Fastest in 1954--63 (1.33 per cent a year), slowest in 1973--83 ($-0.23$). Researchers per million people: United States 1.54 times (1996--2022), Germany 2.11, Japan 1.14,
Korea 4.33, China 4.71 (1996--2023).
\end{solution}

\begin{solution}{exo:b2:endogenous-growth:2}
% code: exercises/solutions.py ak_growth
$0.2 \times 0.5 - 0.05 = 5$ per cent; $0.25 \times 0.5 - 0.05 = 7.5$ per cent.
\end{solution}

\begin{solution}{exo:b2:endogenous-growth:3}
Nonrival: the recipe, the program, the vaccine's formula. Rival: the meal, the computer, the dose.
\end{solution}

\begin{solution}{exo:b2:endogenous-growth:4}
A large, persistent rise in the number of researchers without a matching rise in the growth rate. Yes: the opening chart
and Jones (1995) document it.
\end{solution}

\begin{solution}{exo:b2:endogenous-growth:5}
% code: exercises/solutions.py semi_growth
Higher: $0.01/0.25 = 4$ per cent, against $0.01/0.5 = 2$ per cent: stronger spillovers from existing ideas make growth faster for the same population growth.
\end{solution}

\begin{solution}{exo:b2:endogenous-growth:6}
% code: exercises/solutions.py level_gain
$(3/2)^{1/0.5} \times 0.97/0.98 = 2.25 \times 0.9898 = 2.227$.
\end{solution}

\begin{solution}{exo:b2:endogenous-growth:7}
% code: exercises/solutions.py research_productivity
$g_A/L_A$ changes by $(1/2)/4 = 1/8$: it falls to an eighth.
\end{solution}

\begin{solution}{exo:b2:endogenous-growth:8}
Long-run growth becomes $0.01/0.25 = 4$ per cent, in both scenarios. The level gain becomes $2^{1/0.25} \times 0.96/0.98 = 15.67$,
four times the 3.92 of $\phi = 0.5$. The listing starts both economies on the baseline's balanced path, $a_0 = (0.02(1 - \phi)/n)^{1/(1 - \phi)}$,
which is 0.0625 for $\phi = 0.75$: with $a_0 = 1$ kept, the baseline would not be on its balanced path. The run prints growth
of 8.00 per cent at once, falling to 4.29 after 200 years, and an output ratio of 3.69 after 50 years and 11.84 after 200.
The speed of the transition does not change. Linearised at $z^* = n/(1 - \phi)$, $\dot z = z(n - (1 - \phi)z)$ returns at the rate $n$ whatever $\phi$
is, and the same share of the log-level gain is reached at each date: after 50 years, $\ln(3.692/0.98)/\ln(15.67/0.98) = 0.48$, as
for $\phi = 0.5$, $\ln(1.902/0.98)/\ln(3.918/0.98) = 0.48$.
\end{solution}

\begin{solution}{exo:b2:endogenous-growth:9}
Once a design exists, letting one more firm use it costs nothing, so a competitive price would be zero and the inventor
could not recover the cost of discovery. A patent lets the inventor charge more than marginal cost, which pays for
research but excludes buyers who value the good above its cost of production: a static deadweight loss in exchange for
the incentive to discover.
\end{solution}

\begin{solution}{exo:b2:endogenous-growth:10}
$g_A = \nu s_RLA^{\phi - 1}$. If $g_A$ is constant, the growth rate of the right-hand side is zero: $g_L + (\phi - 1)g_A = 0$, so $g_A = n/(1 - \phi)$.
\end{solution}

\begin{solution}{exo:b2:endogenous-growth:11}
% code: exercises/solutions.py growth_path, time_to
$\dot z/z = g_L + (\phi - 1)g_A = n - (1 - \phi)z$. This logistic equation has solution $z(t) = n/((1 - \phi) + (n/z_0 - (1 - \phi))e^{-nt})$. With $n = 0.01$, $\phi = 0.5$ and $z_0 = 0.04$: $z(t) = 0.01/(0.5 - 0.25e^{-0.01t})$, which equals 0.03 when
$e^{-0.01t} = (0.5 - 1/3)/0.25 = 2/3$, after $t = 100\ln 1.5 = 40.5$ years.
\end{solution}

\begin{solution}{exo:b2:endogenous-growth:12}
From $g_A = \nu s_RLA^{\phi - 1}$: $A^{1-\phi} = \nu s_RL/g_A$. On the balanced path $g_A = n/(1 - \phi)$ does not depend on $s_R$, so for the same $L$, doubling $s_R$ doubles $A^{1-\phi}$ and multiplies $A$ by $2^{1/(1-\phi)}$.
\end{solution}

\begin{solution}{pb:b2:endogenous-growth:1}
% code: exercises/solutions.py brenland; listings/ideas.py
\textbf{Part I.}
1. $\phi = 0.5$: existing ideas help produce new ones, but less than in proportion; ideas get harder to find.
2. $g_A = \nu s_RLA^{-0.5}$.
3. It falls, as $A^{-0.5}$.
4. Each idea makes the next harder; only more researchers keep the flow of ideas growing in proportion to the stock.
5. $g_A/(s_RL) = \nu A^{-0.5}$, falling as $A$ grows.
\textbf{Part II.}
6. $0.01/(1 - 0.5) = 2$ per cent a year.
7. 2 per cent: output per worker is $A(1 - s_R)$ with $s_R$ constant.
8. No.
9. Zero: $n = 0$ gives $g_A = 0$ in the long run.
10. $0.01/0.25 = 4$ per cent.
11. Growth would rise without limit, since $g_A = \nu s_RL$ grows with $L$.
\textbf{Part III.}
12. It falls by about 2 per cent: $0.96/0.98 = 0.980$.
13. It doubles, to 4 per cent.
14. $g_A(t) = 0.01/(0.5 - 0.25e^{-0.01t})$; 3 per cent after $100\ln 1.5 = 40.5$ years.
15. 2.87 per cent.
16. 2 per cent.
17. As ideas accumulate, the same research effort adds a smaller proportion to the stock.
\textbf{Part IV.}
18. $2^{1/0.5} = 4$.
19. $4 \times 0.96/0.98 = 3.918$.
20. After 50 years output per worker is 1.902 times the baseline: in logarithms, $\ln 1.902/\ln 3.918 = 0.47$, about half of the long-run gain.
21. $2^4 \times 0.96/0.98 = 15.67$.
22. Because the level gain is large: almost four times the income per worker in the long run, and nearly twice within fifty
years, at the price of 2 per cent of output at the start.
23. A smaller fall in \omterm{def:b2:endogenous-growth:research-productivity}{research productivity} as ideas accumulate: growth that falls back more slowly after the rise in
research, or long-run growth above 2 per cent with a workforce growing 1 per cent.
24. Brenland's balanced growth rate is $n/(1 - \phi) = 2$ per cent a year whatever its research share; doubling the share leaves it unchanged but multiplies
output per worker in the long run by 3.918.
\end{solution}
